% ============================================================================
%  Mecanismos de consenso en blockchain
%  Módulo 12 · Consenso distribuido — ChainNotes
%
%  Versión LaTeX del apunte web temas/consenso.html
%  Las dos trazas de PoS son las mismas que produce assets/js/consenso.js
%  y las mismas que imprime el documento del curso.
%
%  Compilar:  pdflatex consenso.tex   (dos pasadas, por el índice)
%             o bien:  latexmk -pdf consenso.tex
% ============================================================================

\documentclass[11pt,a4paper]{article}

% --- idioma y codificación -------------------------------------------------
\usepackage[T1]{fontenc}
\usepackage[utf8]{inputenc}
\usepackage[spanish,es-noshorthands,es-tabla]{babel}
\usepackage{lmodern}
\usepackage{microtype}

% --- página ----------------------------------------------------------------
\usepackage[a4paper,margin=2.6cm,headheight=14pt]{geometry}
\usepackage{fancyhdr}
\usepackage{enumitem}
\usepackage{booktabs}
\usepackage{tabularx}
\usepackage{array}
\usepackage{caption}
\usepackage{longtable}

% --- matemáticas -----------------------------------------------------------
\usepackage{amsmath,amssymb,amsthm}
\usepackage{mathtools}

% --- gráficas --------------------------------------------------------------
\usepackage{tikz}
\usetikzlibrary{positioning,arrows.meta,calc,fit,backgrounds,shapes.geometric,
                decorations.pathreplacing}

% --- enlaces y cajas (hyperref antes que tcolorbox) ------------------------
\usepackage[colorlinks=true,linkcolor=violetdeep,urlcolor=violetdeep,
            citecolor=violetdeep,pdftitle={Mecanismos de consenso en blockchain},
            pdfauthor={ChainNotes}]{hyperref}
\usepackage[most]{tcolorbox}

% ---------------------------------------------------------------------------
%  Paleta (heredada de assets/css/styles.css)
% ---------------------------------------------------------------------------
\definecolor{violetc}{HTML}{7A68F5}
\definecolor{violetdeep}{HTML}{5B46E0}
\definecolor{mintc}{HTML}{1FA88A}
\definecolor{coralc}{HTML}{E04E5C}
\definecolor{amberc}{HTML}{C98A1F}
\definecolor{skyc}{HTML}{1E88C7}
\definecolor{inkc}{HTML}{282344}

\hypersetup{linkcolor=violetdeep}

% ---------------------------------------------------------------------------
%  Estilo de párrafo y encabezados
% ---------------------------------------------------------------------------
\setlength{\parindent}{0pt}
\setlength{\parskip}{0.55em}
\captionsetup{font=small,labelfont={bf,color=violetdeep}}

\renewcommand{\topfraction}{0.9}
\renewcommand{\bottomfraction}{0.8}
\renewcommand{\textfraction}{0.07}
\renewcommand{\floatpagefraction}{0.75}
\setcounter{topnumber}{3}
\setcounter{totalnumber}{4}

\pagestyle{fancy}
\fancyhf{}
\fancyhead[L]{\small\color{inkc}\textsc{Mecanismos de consenso}}
\fancyhead[R]{\small\color{inkc}\thepage}
\renewcommand{\headrulewidth}{0.4pt}

% ---------------------------------------------------------------------------
%  Entornos de teoremas
% ---------------------------------------------------------------------------
\theoremstyle{definition}
\newtheorem{definicion}{Definición}
\theoremstyle{plain}
\newtheorem{proposicion}{Proposición}
\theoremstyle{remark}
\newtheorem*{observacion}{Observación}

% ---------------------------------------------------------------------------
%  Cajas de color
% ---------------------------------------------------------------------------
\newtcolorbox{nota}[1]{enhanced,breakable,colback=skyc!6,colframe=skyc!75!black,
  coltitle=white,fonttitle=\bfseries\small,title={#1},
  left=7pt,right=7pt,top=5pt,bottom=5pt,boxrule=0.7pt,arc=2.5pt,
  before upper={\setlength{\parskip}{0.5em}}}

\newtcolorbox{caso}[1]{enhanced,breakable,colback=amberc!8,colframe=amberc!85!black,
  coltitle=white,fonttitle=\bfseries\small,title={#1},
  left=7pt,right=7pt,top=5pt,bottom=5pt,boxrule=0.7pt,arc=2.5pt,
  before upper={\setlength{\parskip}{0.5em}}}

\newtcolorbox{aviso}[1]{enhanced,breakable,colback=coralc!7,colframe=coralc!85!black,
  coltitle=white,fonttitle=\bfseries\small,title={#1},
  left=7pt,right=7pt,top=5pt,bottom=5pt,boxrule=0.7pt,arc=2.5pt,
  before upper={\setlength{\parskip}{0.5em}}}

\newtcolorbox{clave}[1]{enhanced,breakable,colback=mintc!7,colframe=mintc!85!black,
  coltitle=white,fonttitle=\bfseries\small,title={#1},
  left=7pt,right=7pt,top=5pt,bottom=5pt,boxrule=0.7pt,arc=2.5pt,
  before upper={\setlength{\parskip}{0.5em}}}

\newtcolorbox{laboratorio}[1]{enhanced,breakable,colback=violetc!6,colframe=violetc!80!black,
  coltitle=white,fonttitle=\bfseries\small,title={#1},
  left=7pt,right=7pt,top=5pt,bottom=5pt,boxrule=0.7pt,arc=2.5pt,
  before upper={\setlength{\parskip}{0.5em}}}

\newtcolorbox{algoritmo}[1]{enhanced,breakable,colback=inkc!3,colframe=inkc!60,
  coltitle=white,fonttitle=\bfseries\small,title={#1},
  left=9pt,right=7pt,top=5pt,bottom=5pt,boxrule=0.6pt,arc=2.5pt}

% ---------------------------------------------------------------------------
%  Atajos de notación
% ---------------------------------------------------------------------------
\newcommand{\si}{\textcolor{mintc}{\checkmark}}
\newcommand{\no}{\textcolor{coralc}{$\times$}}
\newcommand{\mecanismo}[1]{\textbf{\textcolor{violetdeep}{#1}}}

\setlist[itemize]{label=\textcolor{violetc}{$\bullet$}}

\newlist{opciones}{enumerate}{1}
\setlist[opciones]{label=(\alph*),leftmargin=2.4em,itemsep=0.1em,topsep=0.2em}

% ===========================================================================
\begin{document}
% ===========================================================================

% babel-spanish redefine \% inspeccionando \lastskip, lo que revienta dentro de
% los nodos de TikZ. Restauramos la definición estándar de LaTeX.
\makeatletter
\DeclareRobustCommand{\%}{\char`\%\relax}
\makeatother

\begin{titlepage}
\centering
\vspace*{2.5cm}

{\Large\color{violetdeep}\textsc{ChainNotes}\par}
{\small\color{inkc}Módulo 12 · Consenso distribuido\par}

\vspace{1.4cm}
{\Huge\bfseries Mecanismos de consenso\par}
\vspace{0.4cm}
{\LARGE en blockchain\par}

\vspace{1.2cm}
\begin{minipage}{0.78\textwidth}
\centering\itshape\color{inkc}
Una blockchain es un libro de cuentas compartido que mantienen al mismo tiempo
muchas personas, sin que exista un banco ni una oficina central que lo
administre. Eso plantea una pregunta incómoda: si nadie manda, ¿cómo se ponen de
acuerdo todos sobre cuál es la versión oficial del libro? Hay diez respuestas
serias, y cada una sacrifica algo distinto.
\end{minipage}

\vspace{1.6cm}
\begin{tikzpicture}[scale=1.0]
  \coordinate (S) at (90:2.0);
  \coordinate (V) at (210:2.0);
  \coordinate (D) at (330:2.0);
  \draw[violetc!55,line width=1pt] (S) -- (V) -- (D) -- cycle;
  \foreach \p/\col in {S/mintc,V/amberc,D/skyc}{
    \fill[white] (\p) circle (0.30);
    \draw[\col,line width=0.9pt,fill=\col!18] (\p) circle (0.28);
  }
  \node[mintc,font=\small\bfseries,above=7pt] at (S) {seguridad};
  \node[amberc,font=\small\bfseries,below=8pt] at (V) {velocidad};
  \node[skyc,font=\small\bfseries,below=8pt] at (D) {descentralización};
  \node[violetdeep,font=\small\itshape] at (0,-0.20) {elige dos};
\end{tikzpicture}

\vfill
{\small Apuntes de Blockchain · Proyecto académico\par}
{\small\color{inkc}\today\par}
\end{titlepage}

\tableofcontents
\clearpage

\begin{nota}{Cómo usar este documento}
Es la versión imprimible del apunte interactivo \texttt{temas/consenso.html}.
Cada laboratorio de la página web aparece aquí como un ejemplo resuelto con
números fijos, para que puedas seguir todas las cuentas a mano. Las dos trazas de
Proof of Stake (secciones~\ref{sec:pos} y~\ref{sec:juguete}) son exactamente las
que produce el simulador de la página. Al final hay seis ejercicios con sus
soluciones en el apéndice~\ref{app:soluciones}.

Dos de los mecanismos aquí descritos tienen apunte propio, más detallado:
la prueba de trabajo en \texttt{temas/prueba-de-trabajo.html} y el problema de
los generales bizantinos en \texttt{temas/generales-bizantinos.html}.
\end{nota}

% ===========================================================================
\section{Qué es un mecanismo de consenso}\label{sec:que-es}
% ===========================================================================

Un mecanismo de consenso responde tres preguntas a la vez: \textbf{quién} puede
añadir la siguiente página al libro, \textbf{cómo} se verifica que esa página sea
honesta, y \textbf{qué pasa} con quien intenta hacer trampa.

Cada mecanismo hace una apuesta distinta entre tres cosas que rara vez se
consiguen juntas:

\begin{itemize}
  \item \textbf{Seguridad.} Qué tan caro le sale a un atacante reescribir el
        libro o detenerlo.
  \item \textbf{Velocidad.} Cuántas operaciones por segundo soporta y cuánto
        tarda en confirmarse una.
  \item \textbf{Descentralización.} Cuántas manos distintas deciden, y qué tan
        fácil es entrar a decidir.
\end{itemize}

No es un teorema demostrado como la cota de los generales bizantinos, sino una
observación empírica robusta: los diseños que mejoran mucho en una dimensión
suelen pagarlo en otra.

\subsection{Abierta o con permiso}

Antes de elegir mecanismo hay que saber con quién se juega, y hay dos
configuraciones:

\begin{itemize}
  \item \textbf{Pública y abierta} (sin permiso): cualquiera puede participar y
        leer el libro, como Bitcoin o Ethereum. Nadie conoce a nadie, así que la
        confianza hay que fabricarla.
  \item \textbf{Privada o de consorcio} (con permiso): solo participantes
        conocidos y autorizados, como un grupo de bancos o de empresas. La
        identidad ya está resuelta y eso abarata todo lo demás.
\end{itemize}

\begin{clave}{La regla que explica todo lo que sigue}
Cuanto más abierta y sin confianza previa es la red, más costoso resulta el
consenso. PoW y PoS gastan electricidad o capital precisamente porque no pueden
preguntarle a nadie quién es quién. Cuanto más se conocen los participantes, más
rápido y barato puede ser el acuerdo, como en BFT o PoA, pero a cambio de
depender de que esa lista de conocidos sea honesta.
\end{clave}

% ===========================================================================
\section{Proof of Work}\label{sec:pow}
% ===========================================================================

\begin{nota}{En una imagen}
Imagine un concurso en el que miles de personas intentan resolver un acertijo que
solo se resuelve a fuerza de probar y probar. Quien lo resuelve primero gana el
derecho a escribir la siguiente página, y su esfuerzo es la prueba de que no está
haciendo trampa.
\end{nota}

\subsection{La condición de minado}

Sea $B = (v, h_{\text{prev}}, r_{\text{Merkle}}, t, T, \eta)$ el encabezado del
bloque: versión, hash del bloque anterior, raíz de Merkle de las transacciones,
marca de tiempo, objetivo y \emph{nonce}. Un bloque es válido si
\[
  H(B) \;=\; \mathrm{SHA256}(\mathrm{SHA256}(B)) \;\le\; T,
  \qquad D \;=\; \frac{T_{\max}}{T} \quad\text{(dificultad).}
\]

Como el hash se comporta como un número uniforme en $[0,2^{256})$, cada intento
tiene probabilidad de éxito
\[
  \theta \;=\; \frac{T+1}{2^{256}},
  \qquad
  \mathbb{E}[\text{intentos}] \;=\; \frac{1}{\theta} \;\approx\; \frac{2^{256}}{T}
  \;\approx\; D\cdot 2^{32} \quad\text{(en Bitcoin).}
\]

Si la red calcula $\Lambda$ hashes por segundo y el minero $i$ aporta
$\lambda_i$, el tiempo entre bloques es exponencial y el minero gana en
proporción a su potencia:
\[
  \tau \sim \mathrm{Exp}(\Lambda\theta),
  \qquad \mathbb{E}[\tau] = \frac{1}{\Lambda\theta},
  \qquad \Pr[\text{gana } i] = \frac{\lambda_i}{\Lambda}.
\]

\subsection{Ajuste de dificultad}

Bitcoin recalcula el objetivo cada 2016 bloques, con meta de 600 segundos por
bloque:
\[
  T' \;=\; T \cdot \min\Bigl\{4,\ \max\Bigl\{\tfrac14,\
  \frac{\Delta t_{\text{real}}}{2016\cdot 600\,\mathrm{s}}\Bigr\}\Bigr\}.
\]
El factor se recorta al intervalo $[\frac14, 4]$ para que un salto brusco de
potencia no desestabilice la red.

\subsection{Elección de cadena y seguridad}

La regla de elección de cadena sigue la de mayor \emph{trabajo acumulado}, no la
más larga en número de bloques:
\[
  W \;=\; \sum_j \frac{2^{256}}{T_j + 1}.
\]

\begin{proposicion}[Cota de Nakamoto]
Si un atacante controla una fracción $q < \frac12$ de la potencia y
$p = 1-q$, la probabilidad de que alcance a la cadena honesta desde $z$ bloques
atrás es, en su forma simplificada,
\[
  \Pr[\text{alcanzar desde } z] \;=\; \left(\frac{q}{p}\right)^{z}.
\]
\end{proposicion}

Por eso se esperan varias confirmaciones. Estos son los valores que produce el
laboratorio de la página web:

\begin{center}
\small
\begin{tabular}{@{}rrr@{}}
\toprule
$z$ & $q = 10\%$ & $q = 30\%$ \\
\midrule
1 & $11{,}1\%$      & $42{,}9\%$ \\
3 & $0{,}137\%$     & $7{,}87\%$ \\
6 & $0{,}000\,19\%$ & $0{,}62\%$ \\
9 & $2{,}6\times10^{-7}\%$ & $0{,}047\%$ \\
\bottomrule
\end{tabular}

\vspace{0.4em}
{\footnotesize Probabilidad de que el atacante alcance la cadena honesta.}
\end{center}

\begin{aviso}{Las «seis confirmaciones» no son un teorema}
Son una convención calibrada para un adversario modesto. Con $q = 30\%$ hacen
falta nueve confirmaciones para bajar del $0{,}1\%$. El número seguro depende de
cuánta potencia le supongas al atacante.
\end{aviso}

\begin{algoritmo}{PoW: minar y verificar}
\begin{enumerate}[label=\arabic*:,leftmargin=2.6em,itemsep=0.05em,topsep=0.2em]
  \item \textbf{función} \textsc{Minar}$(B, T, i, n)$
  \item \quad $\eta \leftarrow i$
  \item \quad \textbf{mientras} verdadero \textbf{hacer}
  \item \quad\quad $h \leftarrow \mathrm{SHA256}(\mathrm{SHA256}(B \parallel \eta))$
  \item \quad\quad \textbf{si} $h \le T$ \textbf{entonces} difundir $(B,\eta)$ y devolver
  \item \quad\quad \textbf{si} llega un bloque válido de otro nodo \textbf{entonces}
  \item \quad\quad\quad abandonar; reconstruir $B$ sobre el nuevo bloque
  \item \quad\quad $\eta \leftarrow \eta + n$
  \item \textbf{fin función}
  \item[]
  \item \textbf{función} \textsc{Verificar}$(B,\eta)$ \hfill $\triangleright$ un solo hash
  \item \quad $h \leftarrow \mathrm{SHA256}(\mathrm{SHA256}(B \parallel \eta))$
  \item \quad \textbf{devolver} $h \le T$ \textbf{y} $h_{\text{prev}}$ es la punta conocida
        \textbf{y} las transacciones son válidas
  \item \textbf{fin función}
\end{enumerate}
\end{algoritmo}

\begin{caso}{Dónde conviene}
Pocas escrituras, de alto valor, que nadie debe poder reescribir nunca:
certificación de diamantes, donaciones en contextos de censura, sellado de
títulos y propiedad intelectual. El costo por transacción es irrelevante frente
al valor de lo registrado.
\end{caso}

% ===========================================================================
\section{Proof of Stake}\label{sec:pos}
% ===========================================================================

\begin{nota}{En una imagen}
Es como una asamblea en la que se sortea quién escribe la siguiente página, y
quienes tienen más dinero puesto en garantía tienen más boletos. Si el elegido
hace trampa, pierde parte de su garantía; por eso conviene ser honesto.
\end{nota}

\subsection{Formulación}

Sea $V = \{1,\dots,N\}$ el conjunto de validadores, $b_i$ la garantía
(\emph{stake}) del validador $i$ y $S = \sum_j b_j$. La regla central es
\[
  \Pr[i \text{ propone el bloque del slot } t] = \frac{b_i}{S},
  \qquad
  \mathbb{E}[\text{bloques de } i \text{ en } K \text{ slots}] = K\,\frac{b_i}{S}.
\]

\textbf{Finalidad} (al estilo de Casper FFG): un punto de control $C$ queda
\emph{justificado} si los votos que enlazan un punto justificado $A \to C$ suman
\[
  \sum_{v \,\in\, \text{votos}(A \to C)} b_v \;\ge\; \tfrac{2}{3}\,S,
\]
que en el código se escribe sin fracciones como $3V \ge 2S$.

\textbf{Castigo} (\emph{slashing}): si $v$ firma dos bloques o votos en
conflicto, $b_v \leftarrow (1-\varphi)\,b_v$ y $v$ es expulsado.

\begin{clave}{Tres ideas, y nada más}
\textbf{(i)} un sorteo donde cada validador tiene tantos boletos como monedas en
garantía; \textbf{(ii)} una votación que exige $\frac23$ de la garantía total;
\textbf{(iii)} un pago al proponente si el bloque se acepta. Todo lo demás
—generadores, castigos, expulsiones— son añadidos sobre ese esqueleto.
\end{clave}

\subsection{PoS en su mínima expresión}

Tres validadores, seis boletos, cuatro slots. Los números del sorteo están dados
para que se pueda seguir a mano.

\begin{center}
\small
\begin{tabular}{@{}lrl@{}}
\toprule
Validador & Garantía & Boletos \\
\midrule
Ana   & 3 & $\{0,1,2\}$ \\
Beto  & 2 & $\{3,4\}$ \\
Carla & 1 & $\{5\}$ \\
\midrule
\textbf{Total} & \textbf{6} & \\
\bottomrule
\end{tabular}
\end{center}

Los números del sorteo son $4, 0, 5, 2$. Se desconectan: Beto en el slot 3, y
Beto y Carla en el slot 4. El pago es de 1 moneda al proponente si el bloque se
acepta.

\begin{laboratorio}{Traza de los cuatro slots}
\centering
\small
\begin{tabular}{@{}rrllrrll@{}}
\toprule
$t$ & $r$ & Intervalo & Proponente & Conectados & $V$ & $3V$ vs $2S$ & Resultado \\
\midrule
1 & 4 & $\{3,4\}$   & Beto  & A, B, C & 6 & $18 \ge 12$ & aceptado \\
2 & 0 & $\{0,1,2\}$ & Ana   & A, B, C & 6 & $18 \ge 12$ & aceptado \\
3 & 5 & $\{5\}$     & Carla & A, C    & 4 & $12 \ge 12$ & aceptado \\
4 & 2 & $\{0,1,2\}$ & Ana   & A       & 3 & $9 < 12$    & descartado \\
\bottomrule
\end{tabular}

\vspace{0.6em}
Ganancias finales: Ana 1, Beto 1, Carla 1.
\end{laboratorio}

\textbf{Qué enseña cada slot.}
\begin{itemize}
  \item \textbf{Slot 1:} el caso normal. El número cae en el intervalo de Beto y
        todos votan.
  \item \textbf{Slot 2:} el número más pequeño cae en el intervalo de quien va
        primero en el orden fijo.
  \item \textbf{Slot 3:} Carla, con un solo boleto, también puede ganar. Como
        Beto no vota, $V = 4$ es exactamente $\frac23$ de 6, y eso basta.
  \item \textbf{Slot 4:} Ana propone, pero solo vota ella. Su garantía es la
        mitad del total, que no alcanza los $\frac23$: ni el validador más rico
        puede decidir solo.
\end{itemize}

\begin{aviso}{Proponer no es lo mismo que quedar aceptado}
Ganar el sorteo solo da el \emph{derecho a proponer}. El bloque entra a la cadena
cuando los votos lo aprueban. En el slot 4 Ana gana el sorteo y su bloque se
descarta.
\end{aviso}

\begin{observacion}
El orden en que el algoritmo recorre a los validadores tiene que ser el mismo en
todos los nodos. Si cada uno usara el suyo, elegirían proponentes distintos y no
habría acuerdo: el sorteo debe ser \emph{reproducible}, no sorpresivo. Es la
misma razón por la que el generador es público y determinista.
\end{observacion}

% ===========================================================================
\section{PoS de juguete: ocho slots con todo}\label{sec:juguete}
% ===========================================================================

La versión mínima más tres añadidos: un generador que produce los números del
sorteo, un pago que aumenta la garantía (y con ella los boletos del ganador), y
el castigo por doble firma.

\begin{center}
\small
\begin{tabular}{@{}ll@{}}
\toprule
Parámetro & Valor \\
\midrule
Validadores (en este orden) & Ana, Beto, Carla, Dani \\
Garantías iniciales $b$     & $40,\ 30,\ 20,\ 10 \quad (S = 100)$ \\
Generador                   & $x_t = (21\,x_{t-1} + 7) \bmod 1000,\quad x_0 = 13$ \\
Número del sorteo           & $r_t = x_t \bmod S_t$ \\
Proponente                  & primer $i$ activo con $r_t <$ límite acumulado \\
Finalidad                   & el bloque es final si $3V \ge 2S$ \\
Recompensa                  & $R = 2$ monedas al proponente si el bloque es final \\
Castigo                     & doble firma: $b_p \leftarrow \lfloor b_p/2 \rfloor$ y expulsión \\
Escenario                   & slot 2: Carla firma dos bloques \\
                            & slot 3: Beto desconectado \\
                            & slot 5: Ana desconectada \\
\bottomrule
\end{tabular}
\end{center}

\begin{laboratorio}{Traza esperada de los ocho slots}
\centering
\footnotesize
\begin{tabular}{@{}rrrlrllrll@{}}
\toprule
$t$ & $x_t$ & $S_t$ & Límites (A/B/C/D) & $r_t$ & Proponente & Fuera & $V$ & $3V$ vs $2S$ & $b$ al final \\
\midrule
1 & 280 & 100 & 40/70/90/100 & 80 & Carla & — & 100 & $300 \ge 200$ & 40/30/22/10 \\
2 & 887 & 102 & 40/70/92/102 & 71 & Carla & — & —  & castigo       & 40/30/11$^\dagger$/10 \\
3 & 634 & 80  & 40/70/—/80   & 74 & Dani  & Beto & 50 & $150 < 160$ & 40/30/—/10 \\
4 & 321 & 80  & 40/70/—/80   & 1  & Ana   & — & 80  & $240 \ge 160$ & 42/30/—/10 \\
5 & 748 & 82  & 42/72/—/82   & 10 & Ana   & Ana & — & vacío        & 42/30/—/10 \\
6 & 715 & 82  & 42/72/—/82   & 59 & Beto  & — & 82  & $246 \ge 164$ & 42/32/—/10 \\
7 & 22  & 84  & 42/74/—/84   & 22 & Ana   & — & 84  & $252 \ge 168$ & 44/32/—/10 \\
8 & 469 & 86  & 44/76/—/86   & 39 & Ana   & — & 86  & $258 \ge 172$ & 46/32/—/10 \\
\bottomrule
\end{tabular}

\vspace{0.5em}
\raggedright\footnotesize
$^\dagger$ Carla conserva 11 monedas bloqueadas, pero sale del sorteo: desde el
slot 3 los límites solo incluyen a Ana, Beto y Dani.
\end{laboratorio}

\textbf{Qué se observa.} Los cuatro desenlaces posibles aparecen: bloque final
con recompensa, castigo por doble firma, bloque no final por falta de votos
(50 de 80 es menos de $\frac23$) y slot vacío por proponente desconectado.

\begin{aviso}{El efecto que se le critica al PoS}
Cada recompensa agranda el intervalo del ganador, y con él su probabilidad de
proponer. Ana pasa de 40 a 46 monedas en ocho slots. Es el «el rico se hace más
rico» que aparece en toda discusión sobre participación.
\end{aviso}

\begin{observacion}
En el slot 3 el bloque de Dani es perfectamente válido —las firmas son
correctas, nadie gasta dos veces— y aun así no es final, solo porque faltaron
votos. \textbf{Validez y finalidad son cosas distintas.} Habría bastado con que
Beto estuviera conectado: entonces $V = 80$, $3V = 240 \ge 160$ y el bloque sería
final. Con apenas 30 de garantía Beto no podía proponer nada en ese slot, pero su
voto era decisivo: en PoS la disponibilidad de los validadores es parte de la
seguridad.
\end{observacion}

% ===========================================================================
\section{Participación delegada (DPoS)}\label{sec:dpos}
% ===========================================================================

\begin{nota}{En una imagen}
Es una democracia representativa: los poseedores de monedas no validan
directamente, sino que votan para elegir a un grupo pequeño de representantes que
se turnan para escribir las páginas. Si un representante falla o abusa, puede ser
reemplazado en la siguiente votación.
\end{nota}

Cada poseedor $i$ con participación $s_i$ aprueba un conjunto de candidatos $A_i$
(en EOS, hasta 30). El peso de cada candidato $k$ es
\[
  V_k \;=\; \sum_{i\,:\,k \in A_i} s_i,
  \qquad
  \mathcal{D} \;=\; \mathrm{top}_N\{V_k\}
  \quad (\text{EOS: } N = 21;\ \text{TRON: } N = 27).
\]

Los delegados se turnan en un calendario fijo; si cada uno produce $m$ bloques
seguidos, el productor del slot $t$ es
\[
  d(t) \;=\; \mathcal{D}\bigl[\lfloor t/m \rfloor \bmod N\bigr]
\]
(en el diseño clásico de EOS, $m = 12$ bloques de $0{,}5$\,s: turnos de 6\,s).

\textbf{Irreversibilidad y tolerancia.} Un bloque es irreversible cuando
\[
  \#\{\text{delegados distintos que construyen encima}\}
  \;\ge\; \Bigl\lfloor \tfrac{2N}{3} \Bigr\rfloor + 1
  \quad (15 \text{ de } 21),
\]
\[
  f_{\max} = \Bigl\lfloor \tfrac{N-1}{3} \Bigr\rfloor \quad (6 \text{ de } 21),
\]
donde $f_{\max}$ es el número de delegados caídos o maliciosos que la red tolera.

\textbf{Recompensa típica:} una parte fija por producir y otra según los votos,
\[
  R_k \;=\; \alpha\,\frac{R}{N} \;+\; (1-\alpha)\,R\,\frac{V_k}{\sum_j V_j}.
\]

\begin{observacion}
El castigo aquí es \emph{político}, no una confiscación automática: un delegado
que falla con frecuencia pierde votos y sale de $\mathcal{D}$ en la siguiente
elección. Funciona si los votantes están atentos, que es justo la crítica que se
le hace.
\end{observacion}

% ===========================================================================
\section{Tolerancia a fallos bizantinos (PBFT y variantes)}\label{sec:bft}
% ===========================================================================

\begin{nota}{En una imagen}
Piense en un grupo de generales que deben decidir juntos si atacar o no, sabiendo
que algunos podrían ser traidores. Se envían mensajes y votan en varias rondas
hasta que una mayoría clara coincide; mientras los traidores sean menos de un
tercio, el acuerdo se sostiene.
\end{nota}

Es el mecanismo con más historia formal de todos: su cota de tolerancia está
\emph{demostrada}, no observada. Con $f$ traidores hacen falta al menos
$3f + 1$ participantes:
\[
  n \;>\; 3f.
\]
Por debajo de eso ningún algoritmo funciona, por ingenioso que sea.

Lo que lo distingue en la práctica es la \textbf{finalidad inmediata}: una vez
confirmado, un bloque no se revierte. No hay bifurcaciones temporales ni «esperar
seis confirmaciones». A cambio, la comunicación entre nodos crece
cuadráticamente, y eso es lo que impide llevarlo a miles de participantes.

Sus variantes modernas —Tendermint, HotStuff— reducen ese costo y por eso
aparecen también en redes públicas basadas en participación.

\begin{caso}{Dónde conviene}
Consorcios de pocas decenas de miembros que no confían plenamente entre sí y
necesitan finalidad inmediata: liquidación interbancaria, trazabilidad de tequila
con denominación de origen, consorcio de la industria del diamante.
\end{caso}

% ===========================================================================
\section{Proof of Authority}\label{sec:poa}
% ===========================================================================

\begin{nota}{En una imagen}
Es como un notario: solo pueden firmar las páginas personas o instituciones
previamente autorizadas y con reputación que cuidar. No se compite con dinero ni
con esfuerzo; se confía en la identidad de quien firma.
\end{nota}

Es el extremo opuesto al PoW. Allá la confianza se fabrica quemando electricidad
porque no se sabe quién es quién; aquí se sabe exactamente quién es quién, y lo
que está en juego es su reputación y su responsabilidad legal.

Eso lo vuelve rapidísimo y de costo mínimo, y a la vez lo saca del espíritu
original de no depender de nadie. No es un defecto si el caso de uso lo pide: en
una cadena de frío farmacéutica, que los validadores sean el laboratorio, la
autoridad sanitaria y los hospitales no es una debilidad, es el diseño.

\begin{caso}{Dónde conviene}
Muchas escrituras de bajo valor individual entre entidades ya reguladas: cadena
de frío de vacunas, receta médica electrónica, registro vehicular y de
siniestros. También redes de prueba para desarrolladores.
\end{caso}

% ===========================================================================
\section{Cinco mecanismos con ingenio matemático}\label{sec:modernos}
% ===========================================================================

Los cinco anteriores son la base histórica. Los siguientes resuelven el mismo
problema con una idea matemática distinta.

\subsection{Algorand: sorteo criptográfico secreto}

Imagine una lotería en la que cada participante lleva en el bolsillo un boleto
secreto y solo él sabe si salió premiado. Los ganadores forman un comité que vota
durante unos segundos y después se disuelve. Como nadie sabe de antemano quién
será elegido, un atacante no sabe a quién sobornar o atacar.

El ingenio son las \emph{funciones aleatorias verificables} (VRF): cada quien
calcula un número pseudoaleatorio con su clave secreta y después puede demostrar
a todos que el resultado fue legítimo. Combinado con un protocolo de acuerdo
bizantino, tolera que hasta un tercio de las monedas sean maliciosas.

\subsection{Avalanche: consenso por muestreo aleatorio}

Es como preguntar en una fiesta enorme «¿a qué hora empieza el concierto?». Cada
persona consulta a un puñado de vecinos al azar y, si la mayoría coincide, adopta
esa respuesta y la repite. Tras pocas rondas toda la sala converge, sin que nadie
hable con todos.

Se basa en \emph{muestreo repetido} y en sistemas \emph{metaestables}, como los
cristales que se decantan hacia un estado. Con pocos parámetros —tamaño de la
muestra, umbral de mayoría, número de rondas— se acota la probabilidad de
desacuerdo. El precio: la seguridad es \emph{probabilística}, altísima pero no
certeza absoluta.

\subsection{Prueba de espacio y tiempo (Chia)}

En lugar de gastar electricidad resolviendo acertijos, se reserva espacio en
disco duro: es como llenar una bodega con miles de billetes de lotería ya
preparados. Cada cierto tiempo se saca un número y quien tenga el billete más
cercano gana el turno.

Las pruebas de espacio se apoyan en problemas difíciles de calcular en poco
espacio, y se añade un retardo secuencial verificable para que el tiempo también
cuente. Consume mucha menos energía que el PoW, pero fomenta la compra y el
desgaste de discos duros.

\subsection{Proof of History (Solana)}

Es un reloj de arena criptográfico: una cadena de cálculos que \emph{no se puede
acelerar}, donde cada paso depende del anterior. Al insertar un evento entre dos
pasos se demuestra que ocurrió antes que los siguientes, sin preguntarle la hora
a nadie.

Calcular toma tiempo, pero verificar se hace en paralelo y rápido, así que el
orden de los eventos queda fijado sin comunicación adicional. Por sí sola no es
consenso completo: necesita una votación de tipo BFT encima.

\subsection{Hashgraph: chisme sobre chisme y voto virtual}

Cada nodo cuenta a otro, al azar, todo lo que sabe y \emph{también} quién le
contó qué y cuándo. Con esa red de rumores, cada participante puede deducir por
sí mismo cómo habría votado cada uno de los demás, sin que nadie envíe votos
reales.

Combina \emph{gossip about gossip} con \emph{voto virtual}: el acuerdo se calcula
sobre el grafo de mensajes ya compartidos, sin rondas adicionales de votación. Se
demuestra tolerante a fallos bizantinos en un modelo asíncrono. No es una cadena
de bloques tradicional sino un grafo dirigido acíclico (DAG).

% ===========================================================================
\section{Vista comparativa}\label{sec:comparativa}
% ===========================================================================

\begin{center}
\small
\begin{tabularx}{\textwidth}{@{}lXll@{}}
\toprule
Mecanismo & Quién decide el siguiente bloque & Tipo de red típica & Ejemplo \\
\midrule
PoW   & Quien resuelve primero el acertijo        & Pública, abierta & Bitcoin \\
PoS   & Sorteo según monedas en garantía          & Pública, abierta & Ethereum \\
DPoS  & Representantes elegidos por voto          & Pública, pocos validadores & EOS, TRON \\
BFT   & Votación entre nodos conocidos            & Privada o de consorcio & Cosmos \\
PoA   & Autoridades aprobadas de antemano         & Privada o de pruebas & VeChain \\
\midrule
Algorand   & Comité elegido por sorteo secreto    & Pública, abierta & Algorand \\
Avalanche  & Muestreo aleatorio repetido          & Pública y subredes & Avalanche \\
Espacio y tiempo & Quien aporta más espacio en disco & Pública, abierta & Chia \\
Historia (PoH)   & Líder rotativo con reloj compartido & Pública, participación & Solana \\
Hashgraph  & Voto virtual deducido del grafo      & Validadores conocidos & Hedera \\
\bottomrule
\end{tabularx}
\end{center}

\begin{clave}{Idea para llevarse}
No existe un mecanismo «mejor» en abstracto. Cuanto más abierta y sin confianza
previa es la red, más costoso resulta el consenso (PoW, PoS); cuanto más se
conocen los participantes, más rápido y barato puede ser (BFT, PoA). Elegir un
mecanismo es elegir qué se sacrifica. Los cinco últimos muestran que la matemática
permite mover esas fronteras: sorteos verificables, muestreo aleatorio, relojes
criptográficos y votos virtuales reducen la comunicación o el gasto sin renunciar
del todo a la seguridad.
\end{clave}

% ===========================================================================
\section{Síntesis}\label{sec:sintesis}
% ===========================================================================

\begin{enumerate}[leftmargin=1.5em]
  \item Un mecanismo de consenso responde quién escribe, cómo se verifica y qué
        pasa con el tramposo. Todo lo demás son detalles de implementación.
  \item Seguridad, velocidad y descentralización no se consiguen las tres a la
        vez. Elegir mecanismo es elegir qué sacrificar.
  \item Cuanto más abierta la red, más caro el consenso. La identidad conocida es
        un atajo: BFT y PoA son rápidos porque ya saben quién es quién.
  \item En PoS, ganar el sorteo solo da derecho a \emph{proponer}. El bloque
        entra cuando los votos llegan a $\frac23$ de la garantía.
  \item La seguridad del PoW es física —electricidad quemada— y la del PoS es
        económica —garantía confiscable—. Las dos funcionan; fallan distinto.
  \item Los cinco mecanismos modernos no rompen el trilema: mueven sus fronteras
        con matemáticas, y cada uno paga el movimiento en otra parte.
\end{enumerate}

% ===========================================================================
\section{Ejercicios}\label{sec:ejercicios}
% ===========================================================================

Las soluciones están en el apéndice~\ref{app:soluciones}.

\begin{enumerate}[leftmargin=1.8em,itemsep=0.8em]

\item\label{ej:sorteo} \textbf{El sorteo a mano.}
En la versión mínima, Ana tiene 3 monedas, Beto 2 y Carla 1. Enumera los boletos
de cada uno y di quién propone con $r = 0, 2, 3, 4$ y $5$. Después explica por
qué el orden en que el algoritmo recorre a los validadores tiene que ser el mismo
en todos los nodos.

\item\label{ej:finalidad} \textbf{Validez no es finalidad.}
En la traza del PoS de juguete, el slot 3 produce un bloque de Dani que es
perfectamente válido y aun así no es final. Explica por qué, con los números, y
di qué tendría que haber pasado para que sí lo fuera.

\item\label{ej:confirmaciones} \textbf{Las seis confirmaciones.}
Usando la cota de Nakamoto, calcula el riesgo de que un atacante con el $10\%$ de
la potencia alcance la cadena honesta desde 1, 3 y 6 bloques atrás. ¿Cuántas
confirmaciones harían falta si controlara el $30\%$?

\item\label{ej:dpos} \textbf{La cota de DPoS.}
EOS usa 21 delegados. Calcula $f_{\max}$ y el umbral de irreversibilidad, y
explica por qué un bloque tarda 15 productores distintos en volverse
irreversible y no 11, que sería la mayoría simple.

\item\label{ej:elegir} \textbf{Elegir con requisitos.}
Una universidad quiere registrar los títulos que emite para que cualquiera pueda
verificarlos durante décadas, incluso si la universidad desaparece; el volumen es
bajo. Otra institución quiere registrar las lecturas de 40\,000 sensores de
temperatura, cada minuto. Elige mecanismo para cada una y justifica.

\item\label{ej:rico} \textbf{Discusión.}
El efecto «el rico se hace más rico» del PoS es una crítica frecuente. Formúlala
con precisión usando la traza de la sección~\ref{sec:juguete}, y después
construye el mejor argumento en contra que se te ocurra. ¿Ocurre algo parecido en
PoW?

\end{enumerate}

% ===========================================================================
\appendix
\section{Soluciones}\label{app:soluciones}
% ===========================================================================

\subsection*{Ejercicio \ref{ej:sorteo}}

Los intervalos se forman acumulando garantías en el orden fijo Ana, Beto, Carla:
Ana $\{0,1,2\}$ con límite 3, Beto $\{3,4\}$ con límite 5, Carla $\{5\}$ con
límite 6. Entonces $r = 0$ y $2$ dan \textbf{Ana}; $r = 3$ y $4$ dan
\textbf{Beto}; $r = 5$ da \textbf{Carla}.

El orden importa porque los intervalos dependen de él. Si un nodo recorriera
Carla, Beto, Ana, el boleto 0 sería de Carla y el sorteo daría otro proponente.
Cada nodo calcularía un bloque válido distinto para el mismo slot y no habría
acuerdo.

\subsection*{Ejercicio \ref{ej:finalidad}}

En el slot 3, Carla ya fue expulsada, así que $S = 40+30+10 = 80$. Beto está
desconectado, de modo que solo votan Ana y Dani:
\[
  V = 40 + 10 = 50, \qquad 3V = 150 < 160 = 2S.
\]
Faltan votos para los dos tercios. El bloque de Dani cumple todas las reglas de
validez pero no alcanza la finalidad, así que no cuenta y Dani no cobra.

Habría bastado con que Beto estuviera conectado: entonces $V = 80$,
$3V = 240 \ge 160$ y el bloque sería final.

\subsection*{Ejercicio \ref{ej:confirmaciones}}

Con $q = 0{,}1$ y $p = 0{,}9$, la razón $q/p = 1/9$:
\[
  z=1:\ \tfrac19 \approx 11{,}1\%, \qquad
  z=3:\ \tfrac1{729} \approx 0{,}137\%, \qquad
  z=6:\ \tfrac1{531\,441} \approx 0{,}000\,19\%.
\]

Con $q = 0{,}3$ la razón es $3/7 \approx 0{,}4286$, que cae mucho más despacio.
Para bajar del $0{,}1\%$ hacen falta \textbf{9 confirmaciones}, porque
$0{,}4286^{8} \approx 0{,}0011$ todavía supera el umbral y
$0{,}4286^{9} \approx 0{,}000\,47$ ya no.

El número de confirmaciones «seguro» no es una constante: depende de cuánta
potencia supongas en manos del atacante.

\subsection*{Ejercicio \ref{ej:dpos}}

\[
  f_{\max} = \Bigl\lfloor \tfrac{21-1}{3} \Bigr\rfloor = \lfloor 6{,}67 \rfloor = 6,
  \qquad
  \Bigl\lfloor \tfrac{2\cdot 21}{3} \Bigr\rfloor + 1 = 14 + 1 = 15.
\]

La mayoría simple no basta por la misma razón que en los generales bizantinos.
Si hasta 6 delegados pueden ser maliciosos y respaldar dos historias en
conflicto, dos grupos de 11 podrían compartir hasta 6 tramposos y aun así tener 5
honestos cada uno apoyando cadenas distintas. Con 15 el traslape obliga a que al
menos $15 + 15 - 21 = 9$ delegados estén en ambos grupos, y como los tramposos
son a lo más 6, queda garantizado al menos un honesto en la intersección. Un
honesto no firma dos historias contradictorias, así que no puede haber dos
bloques irreversibles en conflicto.

\subsection*{Ejercicio \ref{ej:elegir}}

\textbf{Los títulos: PoW.} Lo que se compra es permanencia e independencia de la
institución. El volumen es bajo, así que el costo por transacción es
irrelevante, y la propiedad que importa —que nadie pueda reescribir el registro,
ni siquiera la propia universidad dentro de veinte años— es exactamente lo que
PoW ofrece. Basta con anclar el hash de cada título.

\textbf{Los sensores: PoA.} Son 40\,000 escrituras por minuto de bajo valor
individual, entre dispositivos que pertenecen a entidades identificadas. No hay
adversario económico que combatir: se necesita rendimiento alto, costo mínimo y
responsabilidad legal clara de cada validador.

\subsection*{Ejercicio \ref{ej:rico}}

No tiene respuesta única. Elementos que deberían aparecer:

\begin{itemize}
  \item \textbf{La crítica, con números:} Ana pasa de 40 a 46 monedas en ocho
        slots. Su probabilidad de proponer sube de $40/100$ a $46/88$, y cada
        recompensa realimenta la siguiente.
  \item \textbf{Contraargumento 1:} la ganancia es proporcional, no exponencial.
        Todos los validadores honestos ganan en proporción a su garantía, así que
        las \emph{proporciones} relativas apenas se mueven. Aquí Ana sube rápido
        porque la muestra es de ocho slots y porque Carla fue expulsada.
  \item \textbf{Contraargumento 2:} en PoW pasa lo mismo y además con economías
        de escala —hardware más barato por volumen, mejor tarifa eléctrica,
        acceso a financiamiento—. El PoS al menos no tiene descuentos por
        volumen: duplicar la garantía cuesta exactamente el doble.
  \item \textbf{Matiz:} si la recompensa sale de monedas nuevas, los validadores
        ganan frente a los tenedores pasivos, que es una redistribución distinta
        de la que señala la crítica.
\end{itemize}

% ===========================================================================
\section{Referencias}
% ===========================================================================

\begin{itemize}[leftmargin=1.4em]
  \item \emph{Mecanismos de consenso en blockchain. Una primera descripción
        intuitiva para todo público.} Documento del curso. Fuente de las fichas,
        de la formulación de PoW, PoS y DPoS, y de las dos trazas de PoS.
  \item Nakamoto, S. \emph{Bitcoin: A Peer-to-Peer Electronic Cash System}, 2008.
        La sección 11 contiene el cálculo de la probabilidad de alcance.
  \item Buterin, V. y Griffith, V. \emph{Casper the Friendly Finality Gadget},
        2017. Origen de la regla de los dos tercios y del castigo por doble firma.
  \item Lamport, L., Shostak, R. y Pease, M. \emph{The Byzantine Generals
        Problem}, ACM TOPLAS, 1982. La cota $n > 3f$.
  \item Castro, M. y Liskov, B. \emph{Practical Byzantine Fault Tolerance},
        OSDI 1999.
  \item Gilad, Y. et al. \emph{Algorand: Scaling Byzantine Agreements for
        Cryptocurrencies}, SOSP 2017.
  \item Rocket, T. et al. \emph{Scalable and Probabilistic Leaderless BFT
        Consensus through Metastability}, 2019.
  \item Baird, L. \emph{The Swirlds Hashgraph Consensus Algorithm}, 2016.
  \item Temario del Curso de Blockchain, MathCon:
        \url{http://www.math.com.mx/curso_Blockchain.html}
\end{itemize}

\vspace{1em}
\begin{nota}{Apuntes relacionados}
\texttt{temas/prueba-de-trabajo.html} desarrolla PoW con el acertijo minado en
vivo, el ajuste de dificultad y el ataque del 51\%.
\texttt{temas/generales-bizantinos.html} demuestra la cota $n > 3f$ y desarrolla
el algoritmo OM(1).
\end{nota}

\end{document}
